\chapter{\IfLanguageName{dutch}{Stand van zaken}{State of the art}}%
\label{ch:stand-van-zaken}

% Tip: Begin elk hoofdstuk met een paragraaf inleiding die beschrijft hoe
% dit hoofdstuk past binnen het geheel van de bachelorproef. Geef in het
% bijzonder aan wat de link is met het vorige en volgende hoofdstuk.

% Pas na deze inleidende paragraaf komt de eerste sectiehoofding.

Excel wordt doorgaans nog veel gebruikt door bedrijven. Het is een gemakkelijke optie om data op te slaan. Het probleem vormt zich echter wanneer er meerdere bestanden ontstaan. Zo kan data voor een bepaald bestand ergens andere terecht komen. Hierbij komt geen validatie te pas, wat zorgt voor een grote kans op menselijke fouten.
\subsection{De nadelen van Microsoft Excel}
Volgens een studie door \textcite{Aburas2019} zijn er drie types fouten die in een spreadsheet kunnen sluipen. Dit zijn: Mechanical errors, Omission errors en Logic errors. Mechanical errors zijn simpele fouten die gaan over het ingeven van nummers en referenties. Logic errors komen voor wanneer een verkeerde formule gegeven wordt aan een bepaalde cel. Als een formule verkeerd is is dit moeilijk terug te vinden als je niet alles weet van het betreffende vakgebied. Deze twee fouten zijn degene die voorkomen in het geval van Freshfrozen. Weten wat de aard is van deze fouten, en begrijpen wat ze zijn is heel belangrijk om het probleem te kunnen oplossen.
Het concrete probleem met spreadsheets is dat data niet abstract is. de data die je invoert en de berekeningen bevinden zich in hetzelfde bestand \autocite{Aburas2019}. Verder beschrijft \textcite{Fuller2011} dat de intentie om simpele data op te slaan snel kan evolueren in een complex systeem waar fouten snel tevoorschijn komen. Dit volgt dezelfde conclusie van daarnet. De data is niet abstract genoeg, wat zorgt voor complexe bestanden waar snel fouten in kruipen.
\subsection{ERP als oplossing}
Een ERP of Enterprise Resource Planning is een vorm van software gebruikt door veel bedrijven. Een ERP systeem dient als digitale schakel bij bedrijfsprocessen. Dit zorgt ervoor dat deze processen te vereenvoudigen en versnellen over meerdere departementen van een bedrijf heen \autocite{11162534}. Freshfrozen gebruikt op dit moment al een ERP pakket, gemaakt door Mapa. Het stockbeheer van de leeggoedpalletten zit op dit moment nog niet in dit pakket. Een ideale oplossing zou een point-to-point integratie zijn. Dit is een aangepaste implementatie waarbij 2 verschillende producten in direct contact met elkaar staan. dit voldoet zeker aan alle noden van de business, maar is duur, tijdsrovend en moeilijk te onderhouden \autocite{Sah2025}. Als er na deze optie geen mogelijkheid is tot direct implenteren in het pakket zelf is er ook nog een derde optie: Een middleware solution. Dit is een systeem dat zich tussen het eigen systeem en het ERP systeem bevind. Dit reduceert de complexiteit van directe interactie tussen de twee systemen, maar is wel zelf complex om te implementeren en vereist ook continu onderhoud \autocite{Sah2025}. Deze informatie helpt ons op dit moment niet verder om te weten wat de beste optie is. Verder eigen onderzoek is dus vereist.
\subsection{Een eigen applicatie als oplossing}
Dit is de laatste optie. Een aparte applicatie geeft ons complete vrijheid over wat er geïmplementeerd wordt en op welke manier. Volgens een projectrapport door \textcite{MACHADO2023184} is .NET en Blazor WebAssembly een goede optie om een beheerapplicatie te maken. Dit komt door de voordelen op vlak van beveiliging en snelheid. Dit is ook de goedkoopste optie. Volgens een artikel door \textcite{machines13111002} wordt een aangepaste digitale oplossing zelfs aangeraden voor kmo's. zelfs het herwerken van een klein bedrijfsproces kan een katalysator zijn voor digitale maturiteit binnen een bedrijf.

